\documentclass[11pt]{letter}
\usepackage[margin=1in]{geometry}
\usepackage[T1]{fontenc}
\usepackage{mathptmx}
\signature{Stefano Castruccio\\Department of Applied and Computational Mathematics and Statistics\\University of Notre Dame}
\address{Department of Applied and Computational Mathematics and Statistics\\University of Notre Dame\\Notre Dame, IN 46556, USA\\scastruc@nd.edu}
\begin{document}
\begin{letter}{The Editor\\The Annals of Applied Statistics}
\opening{Dear Editor,}

We submit the manuscript ``Escalation or shared episodes? A joint spatio-temporal model of
armed conflict and social unrest'', by Stefano Castruccio and Richard Marcantonio, for
consideration in \emph{The Annals of Applied Statistics}.

Protests and riots are often followed by armed conflict in the same place, and this
association is widely read as escalation and used to forecast conflict. It can also arise
because both kinds of events respond to the same latent episodes, such as a local political
crisis or a period of media attention, with no effect of one on the other. The distinction
matters for the study of political violence and for early warning, and it is a statistical
question: we show that models of conflict alone, including models with their own latent
space-time field, cannot make it, and that a joint model of the two outcomes can.

The paper proposes a joint latent Gaussian model for the weekly occurrence of conflict and
unrest on a $0.5^\circ$ grid, in which the two outcomes share a space-time field with an
estimated loading, alongside cross-lagged effects of their recent histories. The fields are
discretised by finite volumes directly on the sphere, with sea barriers, and fitted by an
empirical Bayes Laplace approximation; case-control subsampling of non-events makes the
hyperparameter estimation feasible on millions of cell-weeks per outcome. Simulation studies
show that univariate models report escalation that does not exist, while the joint model has
close to nominal coverage. In ACLED data for West Africa and the Sahel and for Africa and the
Middle East, 2018--2024, shared episodes account for between a quarter and two thirds of the
association between recent unrest and conflict. We also examine what the joint model adds to
prediction, and find that its value lies in filling in an outcome where it goes unrecorded
while the other is observed, rather than in forecasting when both are observed.

We believe the paper suits the journal because the substantive question drives the
statistical development, and because the methods, a multivariate latent Gaussian model on
the sphere with scalable inference for rare binary outcomes, apply to other multivariate
event data in the social and environmental sciences. The Supplementary Material contains the
construction of the data, computational details, simulation studies of the univariate
building blocks, and further results. Code for all analyses will be made available; the
ACLED data are available from ACLED under its terms of use.

The manuscript is not under consideration elsewhere, and both authors have approved the
submission. The use of a generative AI assistant in preparing the code and the text is
described in the Acknowledgments. We have no competing interests to declare.

\closing{Sincerely,}
\end{letter}
\end{document}
